\chapter*{Abstract}
\addcontentsline{toc}{chapter}{Abstract}

The \textbf{Bank Loan Management System} is a comprehensive software engineering specification and architectural design formulated to modernize and automate the credit management lifecycle within retail commercial banking institutions. Conventional paper-based lending procedures suffer from severe structural liabilities, including prolonged processing lead times, repetitive physical documentation, human underwriting errors, lack of applicant visibility, and high tracking overhead.

This project delivers an end-to-end software engineering documentation blueprint adhering to IEEE standards and university curriculum requirements (CIC-357). The project encompasses rigorous requirement engineering culminating in a detailed Software Requirement Specification (SRS) defining 18 functional requirements and 8 non-functional quality attributes. System modeling is executed through two complementary engineering paradigms: function-oriented modeling (Data Flow Diagrams Levels 0, 1, 2, and Structure Chart) and object-oriented modeling across the complete suite of Unified Modeling Language (UML) diagrams (Use Case, Class, Object, State-Chart, Activity, Sequence, Collaboration, Component, and Deployment diagrams).

The database architecture is designed in Third Normal Form (3NF) across six core relational entities supported by an exhaustive data dictionary and integrity constraints. Software verification is documented via 22 black-box test cases spanning Equivalence Partitioning and Boundary Value Analysis. Quantitative project metrics are derived through Albrecht / IFPUG Function Point Analysis (191.4 Adjusted Function Points, ~9,570 SLOC, 37.5 Person-Months) and network scheduling via CPM/PERT identifying a 99-day critical path. The resulting documentation establishes a sound, publication-grade foundation bridging academic software engineering theory with modern enterprise FinTech architecture.
